%! Piecewise Linear Learning Functions — Unit 8, Lesson 8.1 — Guided Notes (Answer Key)
\documentclass[10pt]{article}
\usepackage{linalg-article}
\usepackage{linalg-key}   % pulls in -boxes; do NOT also load -boxes
\newcommand{\TallMath}[1]{$\displaystyle #1\rule[-1.4em]{0pt}{3.2em}$}
\newcommand{\vv}{\mathbf{v}}
\newcommand{\bb}{\mathbf{b}}
\newcommand{\zero}{\mathbf{0}}
\newcommand{\relu}{\mathrm{ReLU}}

\begin{document}

\pageheader{Unit 8, Lesson 8.1}{Guided Notes --- Answer Key}
\namedateperiod

\vspace{0.12in}

\begin{objectivebox}
\textbf{By the end of this lesson, I will be able to\dots}
\begin{itemize}\setlength{\itemsep}{2pt}
  \item build a \emph{continuous piecewise-linear} function by shifting and adding ReLUs, and read off where each \emph{fold} sits;
  \item find the slope on each straight piece as the running sum of the switched-on ReLUs, and count pieces: $N$ folds make $N+1$ pieces;
  \item explain how a layer $\relu(A\vv+\bb)$ produces these bending functions --- the bias places each fold, the weights set each slope --- so a wider layer can fit wigglier data.
\end{itemize}
\end{objectivebox}

\vspace{0.06in}

\begin{vocabbox}
\small
Fill in each term as we build it together.
\vspace{2pt}
\vocabans{Piecewise-linear function (straight pieces joined at folds)}{a function made of straight segments meeting at corners; a sum of shifted ReLUs is one, and it can trace any bending data}
\vocabans{Fold (corner) --- where the slope changes}{a point where the graph bends; each ReLU contributes exactly one, at its shift}
\vocabans{Shift $\relu(x-s)$ --- moves a fold to $x=s$ (the bias)}{replacing $x$ with $x-s$ slides the fold from $0$ to $s$; inside a layer the bias $\bb$ does this}
\vocabans{Weight on a ReLU --- how much slope it adds}{the coefficient $c$ in $c\,\relu(x-s)$ sets how much the slope changes at that fold; $c<0$ bends the graph down}
\vocabans{Counting rule: $N$ folds $\Rightarrow N+1$ straight pieces}{each new fold starts a new straight piece, so $N$ ReLUs make a graph with $N+1$ linear pieces}
\end{vocabbox}

\begin{hookbox}
Last lesson one \textbf{ReLU} gave us a single \emph{fold}, and one \textbf{layer} $\relu(A\vv+\bb)$ ran a
matrix step and then bent. Today we look at the \emph{shape} that bending makes.
\begin{itemize}\setlength{\itemsep}{1pt}
  \item A single straight line $y=wx+b$ has one slope forever. To trace data that rises \emph{and} falls, what does the graph need that a line does not have? \ansline{It needs \emph{folds} (corners) where the slope can change --- so it can go up on one piece and down on the next.}
  \item If replacing $x$ with $x-2$ inside $\relu$ slides its fold from $x=0$ to $x=2$, what part of a layer ($A\vv+\bb$) do you think controls \emph{where} a fold sits? \ansline{The \emph{bias} $\bb$ --- adding a constant inside the ReLU shifts the fold, exactly like $x-2$.}
\end{itemize}
\end{hookbox}

\begin{notesbox}{1. One Fold Is Not Enough --- We Add Folds}
A single ReLU bends once. Real data bends many times, so we build a \textbf{piecewise-linear} function ---
several straight pieces joined at corners --- by \emph{adding shifted ReLUs together}. Two knobs control the
build:
\begin{itemize}\setlength{\itemsep}{2pt}
  \item \textbf{Shift $=$ where.} $\relu(x-s)$ is flat until $x=s$, then rises: its fold sits at $x=s$. The
        \emph{bias} inside a layer is this shift --- it \emph{places} each fold.
  \item \textbf{Weight $=$ how steep.} A coefficient $c$ in front, $c\cdot\relu(x-s)$, decides how much the
        slope changes at that fold (and $c<0$ bends the graph \emph{down}).
\end{itemize}
Place folds with shifts, set slopes with weights, and a sum of ReLUs can trace any bending data.
\end{notesbox}

\begin{notesbox}{2. Building a Function: the Tent}
Let us add three ReLUs into one function:
\[
  F(x)=\relu(x)\;-\;2\,\relu(x-2)\;+\;\relu(x-4).
\]
Each piece switches on at its fold: the first at $x=0$, the second at $x=2$, the third at $x=4$. Evaluate
$F$ by turning on only the ReLUs whose fold is to the \emph{left} of $x$:
\[
  F(0)={\color{keyred}\mathbf{0}},\qquad F(2)={\color{keyred}\mathbf{2}},\qquad F(4)={\color{keyred}\mathbf{0}},\qquad F(6)={\color{keyred}\mathbf{0}}.
\]
(For $F(2)$: only $\relu(x)$ is on, so $F(2)=2$. For $F(4)$: $\relu(x)=4$ and $-2\,\relu(x-2)=-2(2)=-4$, so
$F(4)=4-4=0$.) The graph rises to a peak, comes back down, then stays flat: a \textbf{tent}. Three folds,
built from three ReLUs.

\begin{center}
\begin{tikzpicture}[scale=0.62]
  % LEFT: a shift slides the fold
  \begin{scope}[shift={(-3.2,0)}]
    \draw[->, charcoal, thin] (-0.4,0)--(4.4,0) node[right]{\scriptsize $x$};
    \draw[->, charcoal, thin] (0,-0.5)--(0,2.9) node[above]{\scriptsize $y$};
    \draw[ultra thick, royalblue] (0,0)--(2,2) node[right]{\scriptsize $\relu(x)$};
    \draw[ultra thick, burgundy] (0,0)--(2,0)--(3.7,1.7);
    \node[burgundy] at (3.3,0.55) {\scriptsize $\relu(x-2)$};
    \fill[burgundy] (2,0) circle (1.6pt);
    \node[charcoal] at (2.0,-0.9) {\scriptsize the shift slides the fold};
  \end{scope}
  % ARROW
  \draw[->, very thick, goldacc] (1.3,1.0)--(3.2,1.0) node[midway, above]{\scriptsize add};
  % RIGHT: the assembled tent
  \begin{scope}[shift={(5.0,0)}]
    \draw[->, charcoal, thin] (-0.4,0)--(5.4,0) node[right]{\scriptsize $x$};
    \draw[->, charcoal, thin] (0,-0.5)--(0,2.9) node[above]{\scriptsize $F(x)$};
    \draw[ultra thick, burgundy] (-0.3,0)--(0,0)--(2,2)--(4,0)--(5.2,0);
    \fill[burgundy] (0,0) circle (1.6pt);
    \fill[burgundy] (2,2) circle (1.6pt);
    \fill[burgundy] (4,0) circle (1.6pt);
    \node[charcoal] at (2.05,2.42) {\scriptsize peak $(2,2)$};
    \node[charcoal] at (0.0,-0.9) {\scriptsize fold};
    \node[charcoal] at (2.0,-0.9) {\scriptsize fold};
    \node[charcoal] at (4.0,-0.9) {\scriptsize fold};
  \end{scope}
\end{tikzpicture}
\end{center}
\end{notesbox}

\begin{notesbox}{3. Reading Slopes and Counting Pieces}
On each straight piece, the slope is the \textbf{sum of the weights of the ReLUs that have switched on} so
far (folds to the left). Walk the tent $F(x)=\relu(x)-2\,\relu(x-2)+\relu(x-4)$ left to right:
\[
  \underbrace{x<0}_{\text{none on}}:\ \text{slope}=0,\qquad
  \underbrace{0<x<2}_{+\relu(x)}:\ \text{slope}={\color{keyred}\mathbf{1}}.
\]
\[
  \underbrace{2<x<4}_{-2\,\relu(x-2)}:\ \text{slope}=1-2={\color{keyred}\mathbf{-1}},\qquad
  \underbrace{x>4}_{+\relu(x-4)}:\ \text{slope}={\color{keyred}\mathbf{0}}.
\]
Each fold \emph{adds} its weight to the slope: $+1$ at $x=0$, then $-2$ at $x=2$, then $+1$ at $x=4$. Notice
the shape: up, then down, then flat --- four straight pieces.

\vspace{2pt}
\textbf{Counting rule.} Three ReLUs gave three folds and \ans{$4$} straight pieces. In general,
\[
  N \text{ folds} \;\Longrightarrow\; N+1 \text{ straight pieces.}
\]
More ReLUs $\Rightarrow$ more folds $\Rightarrow$ more pieces $\Rightarrow$ a function that can bend more.
\end{notesbox}

\begin{notesbox}{4. This Is Exactly a Layer --- and Why It Matters}
Where do all these shifted ReLUs come from? From \textbf{one layer} $\relu(A\vv+\bb)$. Feed in a single
number $x$; a weight matrix with several rows makes several copies $a_i x + b_i$, and the ReLU bends each
one. Row $i$'s bias $b_i$ \emph{places} fold number $i$, and its weight $a_i$ \emph{sets the slope} there.
So a layer with $N$ rows (ReLUs) builds a piecewise-linear function with \ans{$N$} folds --- one per row.

\vspace{2pt}
\textbf{Why it matters.} This is how a network fits data: pick the weights and biases so the folds land in
the right places and the pieces have the right slopes, and the graph threads the data points. A \emph{wider}
layer (a bigger matrix $A$) means more folds --- so it can trace \emph{wigglier} data. Choosing those
weights to shrink the loss is \textbf{learning} (Lesson 8.3).

\vspace{2pt}
\textbf{The road ahead.} \textbf{8.2} convolutional neural nets\; $\bullet$\; \textbf{8.3} minimizing loss
by gradient descent\; $\bullet$\; \textbf{8.4} mean, variance, and covariance.
\end{notesbox}

\begin{practicebox}
Show your work.
\begin{enumerate}[leftmargin=1.6em, itemsep=4pt]
  \item For $G(x)=\relu(x)-\relu(x-2)$: where are the folds, what is the slope on each of the three pieces, and what are $G(1)$ and $G(3)$? \ansline{Folds at $x=0$ and $x=2$. Slopes $0$ (for $x<0$), $1$ (for $0<x<2$), $0$ (for $x>2$). $G(1)=1$ and $G(3)=3-1=2$ --- a ramp that levels off at height $2$.}
  \item A piecewise-linear function is built from $4$ ReLUs. How many folds does it have, and how many straight pieces? \ansline{$4$ folds and $4+1=5$ straight pieces.}
  \item For the tent $F(x)=\relu(x)-2\,\relu(x-2)+\relu(x-4)$, what is $F(2)$, and what is the slope of the piece just after $x=2$? \ansline{$F(2)=2$ (the peak). Just after $x=2$ the $-2\,\relu(x-2)$ switches on, so the slope is $1-2=-1$ (the graph turns downward).}
\end{enumerate}
\end{practicebox}

\begin{teachernote}
The spine is Section~2 (build the tent $F=\relu(x)-2\,\relu(x-2)+\relu(x-4)$) and Section~3 (read slopes as
a \emph{running sum of switched-on weights}, and count $N$ folds $\to N+1$ pieces). Section~1 gives the two
knobs --- \emph{shift places the fold, weight sets the slope} --- and Section~4 lands the payoff: these
shifted ReLUs are just the rows of a layer $\relu(A\vv+\bb)$, so bias $=$ fold location and width $=$ number
of folds. Keep the arithmetic light and lean on the picture. Most common slips: forgetting that a ReLU only
turns on once $x$ passes its fold (so at a given $x$ you sum only the folds to the \emph{left}); adding a
ReLU's \emph{value} to the slope instead of its \emph{weight}; and miscounting pieces (it is folds $+1$).
Sentence to land: shift to place each fold, weight to set each slope --- that is how a layer bends to fit
data.
\end{teachernote}

\end{document}
